\section*{Chapter \ref{ch:b1:quantum-numbers} --- Quantum Numbers and Electron Configurations}

\begin{solution}{exo:b1:quantum-numbers:1}
(a) Impossible: $l \le n - 1 = 1$. (b) Possible (a 3p electron).
(c) Impossible: $m_s = \pm\tfrac12$. (d) Possible (a 4f electron).
(e) Impossible: for $l = 0$, $m_l$ can only be 0.
\end{solution}

\begin{solution}{exo:b1:quantum-numbers:2}
3d: 5 orbitals, 10 electrons; 4f: 7 orbitals, 14 electrons; 5p: 3
orbitals, 6 electrons; shell $n = 2$: $n^2 = 4$ orbitals, $2n^2 = 8$
electrons.
\end{solution}

\begin{solution}{exo:b1:quantum-numbers:3}
\ce{O}: $1\text{s}^2 2\text{s}^2 2\text{p}^4 = [\ce{He}]\,2\text{s}^2 2\text{p}^4$,
6 \omterm{def:b1:quantum-numbers:configuration}{valence electrons}. \ce{P}: $1\text{s}^2 2\text{s}^2 2\text{p}^6 3\text{s}^2
3\text{p}^3 = [\ce{Ne}]\,3\text{s}^2 3\text{p}^3$, 5. \ce{K}: $[\ce{Ar}]\,4\text{s}^1$,
1. \ce{Ca}: $[\ce{Ar}]\,4\text{s}^2$, 2. \ce{Br}: $[\ce{Ar}]\,3\text{d}^{10}
4\text{s}^2 4\text{p}^5$, 7 (the filled 3d belongs to the core).
% ledger: gs:K, gs:Ca
\end{solution}

\begin{solution}{exo:b1:quantum-numbers:4}
\ce{N}, 2p: \omorbs{u,u,u}, 3 unpaired. \ce{S}, 3p: \omorbs{ud,u,u}, 2
unpaired. \ce{Cl}, 3p: \omorbs{ud,ud,u}, 1 unpaired. In each case 2s or
3s is full, \omorbs{ud}.
\end{solution}

\begin{solution}{exo:b1:quantum-numbers:5}
\ce{Na+}, \ce{Mg^{2+}}, \ce{Al^{3+}}, \ce{F-}, \ce{O^{2-}}: $1\text{s}^2
2\text{s}^2 2\text{p}^6$, isoelectronic with neon. \ce{Cl-}, \ce{S^{2-}},
\ce{Ca^{2+}}: $[\ce{Ne}]\,3\text{s}^2 3\text{p}^6$, isoelectronic with
argon.
\end{solution}

\begin{solution}{exo:b1:quantum-numbers:6}
The 4s electrons leave first (\cref{prop:b1:quantum-numbers:ions}).
\ce{Ti^{2+}}: $[\ce{Ar}]\,3\text{d}^2$, 2 unpaired; \ce{Mn^{2+}}:
$3\text{d}^5$, 5; \ce{Co^{2+}}: $3\text{d}^7$, 3 (\omorbs{ud,ud,u,u,u});
\ce{Ni^{2+}}: $3\text{d}^8$, 2; \ce{Zn^{2+}}: $3\text{d}^{10}$, 0.
% ledger: gs:Ti2+, gs:Mn2+, gs:Co2+, gs:Zn2+
\end{solution}

\begin{solution}{exo:b1:quantum-numbers:7}
$\Delta E = 13.6\,(1/9 - 1/16) = 13.6 \times 7/144 = \qty{0.661}{eV}$;
$\lambda = 1240/0.661 = \qty{1876}{nm} \approx \qty{1.88}{\micro m}$:
infrared (first line of the Paschen series).
\end{solution}

\begin{solution}{exo:b1:quantum-numbers:8}
(a) $[\ce{Ne}]\,3\text{s}^2 3\text{p}^4$: sulfur, $Z = 16$.
(b) $[\ce{Ar}]\,3\text{d}^{10} 4\text{s}^2 4\text{p}^5$: bromine, $Z = 35$.
(c) $[\ce{Kr}]\,5\text{s}^1$: rubidium, $Z = 37$.
(d) $[\ce{Ar}]\,3\text{d}^3 4\text{s}^2$: vanadium, $Z = 23$.
A 3p electron of sulfur: $(3, 1, -1, +\tfrac12)$, for instance.
% ledger: gs:V
\end{solution}

\begin{solution}{exo:b1:quantum-numbers:9}
For 4s, $n + l = 4 + 0 = 4$; for 3d, $n + l = 3 + 2 = 5$: 4s comes
first. After 3p ($n + l = 4$) the next \omterm{def:b1:quantum-numbers:orbital}{subshell} is 4s ($n + l = 4$, larger
$n$), which opens period 4; 3d ($n + l = 5$) comes only after 4s.
Period 3 therefore contains 3s and 3p only: $2 + 6 = 8$ elements.
\end{solution}

\begin{solution}{exo:b1:quantum-numbers:10}
Ionisation energy $Z^2E_R = 4 \times 13.6 = \qty{54.4}{eV}$, four times
that of hydrogen. $2 \to 1$: $\Delta E = 54.4\,(1 - 1/4) = \qty{40.8}{eV}$,
$\lambda = 1240/40.8 = \qty{30.4}{nm}$, four times shorter than the
hydrogen line at \qty{121.6}{nm}: every energy is multiplied by $Z^2 = 4$.
\end{solution}

\begin{solution}{exo:b1:quantum-numbers:11}
(a) \omterm{def:b1:quantum-numbers:energy-level}{Excited state} of beryllium (\omterm{def:b1:quantum-numbers:energy-level}{ground state} $1\text{s}^2 2\text{s}^2$).
(b) Impossible: 2p holds at most 6 electrons. (c) \omterm{def:b1:quantum-numbers:energy-level}{Excited state} of
beryllium. (d) \omterm{def:b1:quantum-numbers:energy-level}{Ground state} of phosphorus. (e) Impossible: an s \omterm{def:b1:quantum-numbers:orbital}{subshell}
holds two electrons (Pauli).
\end{solution}

\begin{solution}{exo:b1:quantum-numbers:12}
\ce{Fe^{3+}} $3\text{d}^5$: 5; \ce{Mn^{2+}} $3\text{d}^5$: 5; \ce{Fe^{2+}}
$3\text{d}^6$: 4; \ce{Cu^{2+}} $3\text{d}^9$: 1; \ce{Zn^{2+}}
$3\text{d}^{10}$: 0. Ranking: $\ce{Fe^{3+}} = \ce{Mn^{2+}} > \ce{Fe^{2+}}
> \ce{Cu^{2+}} > \ce{Zn^{2+}}$. Iron(III) has one more \omterm{def:b1:quantum-numbers:unpaired}{unpaired electron}
than iron(II), so it is attracted more; zinc(II) has none, so it is not
attracted.
% ledger: gs:Fe2+, gs:Fe3+, gs:Mn2+, gs:Cu2+, gs:Zn2+
\end{solution}

\begin{solution}{pb:b1:quantum-numbers:1}
\textbf{1.} $E_1 = \qty{-13.6}{eV}$, $E_2 = \qty{-3.40}{eV}$,
$E_3 = \qty{-1.51}{eV}$, $E_4 = \qty{-0.85}{eV}$.
\textbf{2.} $\Delta E = 3.40 - 1.51 = \qty{1.89}{eV}$, $\lambda = 1240/1.89
= \qty{656}{nm}$: red.
\textbf{3.} $4 \to 2$: \qty{2.55}{eV}, \qty{486}{nm} (blue-green);
$5 \to 2$: \qty{2.86}{eV}, \qty{434}{nm} (violet); $6 \to 2$: \qty{3.022}{eV},
\qty{410}{nm} (violet).
\textbf{4.} $7 \to 2$: \qty{3.12}{eV}, \qty{397}{nm}, already ultraviolet;
the higher transitions are shorter still, and $3 \to 2$ is the longest
Balmer wavelength. The Lyman lines are all ultraviolet and the Paschen
lines all infrared, so only four lines are visible.
\textbf{5.} $\infty \to 2$: \qty{3.40}{eV}, $\lambda = \qty{365}{nm}$.
\textbf{6.} $\Delta E = 13.6 \times 3/4 = \qty{10.2}{eV}$, $\lambda =
\qty{122}{nm}$: far ultraviolet, invisible to the eye and absorbed by
glass and air.
\textbf{7.} The photon must bring at least \qty{13.6}{eV}: $\lambda \le
1240/13.6 = \qty{91.2}{nm}$.
\textbf{8.} $l = 0$ ($m_l = 0$), $l = 1$ ($m_l = -1, 0, 1$), $l = 2$
($m_l = -2, \dots, 2$): $1 + 3 + 5 = 9$ orbitals.
\textbf{9.} $\sum_{l=0}^{n-1}(2l+1) = n^2$ orbitals of two electrons:
$2n^2$; for $n = 4$, 32.
\textbf{10.} 1s, 2s, 2p, 3s, 3p, 4s, 3d, 4p, 5s, 4d, 5p, 6s, 4f, 5d, 6p,
7s, 5f, 6d, 7p.
\textbf{11.} Each period runs from $n$s to $n$p: 2, 8, 8, 18, 18, 32, 32.
\textbf{12.} 3d ($n + l = 5$) comes after 4s ($n + l = 4$), which opens
period 4.
\textbf{13.} $Z = 2, 10, 18, 36, 54, 86, 118$.
\textbf{14.} $[\ce{Ar}]\,3\text{d}^6 4\text{s}^2$; 3d: \omorbs{ud,u,u,u,u},
4 \omterm{def:b1:quantum-numbers:unpaired}{unpaired electrons}.
\textbf{15.} \ce{Fe^{2+}} $[\ce{Ar}]\,3\text{d}^6$: 4 unpaired;
\ce{Fe^{3+}} $[\ce{Ar}]\,3\text{d}^5$: 5 unpaired, a half-filled \omterm{def:b1:quantum-numbers:orbital}{subshell}.
\textbf{16.} Predicted $3\text{d}^4 4\text{s}^2$: 4 unpaired; observed
$3\text{d}^5 4\text{s}^1$: 6 unpaired (5 in 3d, 1 in 4s).
\textbf{17.} \ce{Cu} $[\ce{Ar}]\,3\text{d}^{10} 4\text{s}^1$, \ce{Cu+}
$[\ce{Ar}]\,3\text{d}^{10}$, \ce{Cu^{2+}} $[\ce{Ar}]\,3\text{d}^9$.
% ledger: gs:Cu, gs:Cu+, gs:Cu2+
\textbf{18.} \ce{Cu^{2+}}, with one \omterm{def:b1:quantum-numbers:unpaired}{unpaired electron}; \ce{Cu+} has none.
\textbf{19.} \ce{Sc^{3+}} $[\ce{Ar}]$: 0; \ce{Ti^{3+}} $[\ce{Ar}]\,3\text{d}^1$:
1; \ce{Mn^{2+}} $[\ce{Ar}]\,3\text{d}^5$: 5. \ce{Mn^{2+}} has the most.
\textbf{20.} Three electrons per orbital; $3(2l + 1)$ per \omterm{def:b1:quantum-numbers:orbital}{subshell};
$3n^2$ per shell.
\textbf{21.} 1s: 3; 2s: 3; 2p: 9; 3s: 3; 3p: 9.
\textbf{22.} The first period closes when 1s is full: $Z = 3$.
\textbf{23.} $Z = 4$, one electron in 2s beyond the noble core.
\textbf{24.} 2s and 2p: $3 + 9 = 12$ elements.
\textbf{25.} $3 + 12 = 15$: the second noble gas of the three-spin world
has \textbf{$Z = 15$}.
\end{solution}
